%==============================================
% DISCUSSION
%==============================================
\section{Discussion}

\subsection{Key Findings}

\begin{enumerate}
\item \textbf{Automated BAP generation is feasible}: Our pipeline successfully automates the tedious manual process of BAP specification, reducing per-structure processing time from 30--60 minutes to $<$1 second.

\item \textbf{High stereochemistry accuracy}: With $>$98\% stereochemistry accuracy, GlycoSMILES2BAP approaches the quality of manual specification while being fully automated.

\item \textbf{Practical usability}: The tool integrates with existing glycan databases (GlyTouCan, GlyGen) and accepts standard notations (IUPAC, WURCS), making it immediately useful for the glycobiology community.

\item \textbf{Each module contributes significantly}: Ablation studies (Table~\ref{tab:ablation}) confirm that all three core modules---CCD Mapper, Anomeric Tracker, and Branch Handler---are essential for high accuracy.
\end{enumerate}

\subsection{Strengths}

\begin{itemize}
\item \textbf{Automation}: Eliminates the need for manual atom-by-atom bond specification, democratizing AF3 glycan modeling for non-experts.

\item \textbf{Accuracy}: Achieves $>$98\% stereochemistry accuracy, approaching the theoretical maximum of manual specification.

\item \textbf{Speed}: Sub-second processing time enables high-throughput applications and database-scale predictions.

\item \textbf{Extensibility}: The modular architecture allows easy addition of new monosaccharide CCD mappings as needed.

\item \textbf{Open source}: Freely available implementation facilitates community contribution and validation.
\end{itemize}

\subsection{Limitations}

\begin{enumerate}
\item \textbf{CCD coverage}: While 28+ monosaccharides are supported, some rare or modified sugars require custom CCD templates not currently available in the PDB database. Specifically, GlcN (glucosamine) and GalN (galactosamine) lack standard CCD entries.

\item \textbf{Input format dependency}: The pipeline relies on correct IUPAC/WURCS notation; errors in input strings propagate through the system.

\item \textbf{Validation scope}: Our benchmark focuses on common mammalian glycans; bacterial and plant glycans with unusual linkages may require additional testing.

\item \textbf{AF3 dependency}: The tool is designed specifically for AF3 input format and may require modification for other structure prediction tools.

\item \textbf{No structural validation}: The pipeline generates input specifications but does not validate the final AF3 predictions against experimental structures.
\end{enumerate}

\subsection{Comparison with Existing Tools}

\begin{table}[h]
\centering
\caption{Comparison of glycan structure prediction input methods.}
\label{tab:comparison}
\begin{tabular}{llccc}
\toprule
\textbf{Tool} & \textbf{Input Format} & \textbf{AF3 Compatible} & \textbf{Automation} & \textbf{Accuracy} \\
\midrule
GlycoSMILES2BAP & IUPAC/WURCS & Yes (CCD+BAP) & Full & $>$98\% \\
Manual BAP & Custom & Yes & None & $\sim$100\% \\
SMILES input & SMILES & Yes & Full & $\sim$60\% \\
userCCD & Custom CCD & Yes & Partial & $\sim$80\% \\
GlyLES & IUPAC & No & Full & N/A \\
\bottomrule
\end{tabular}
\end{table}

\subsection{Community Impact}

GlycoSMILES2BAP addresses a critical bottleneck identified by \citet{huang2025}, enabling:

\begin{enumerate}
\item \textbf{Large-scale screening}: Researchers can now predict structures for hundreds of glycan variants without manual specification.

\item \textbf{Reproducibility}: Automated conversion eliminates human error in BAP specification.

\item \textbf{Integration}: Direct compatibility with GlyTouCan (200,000+ glycan entries) enables systematic structural annotation.

\item \textbf{Education}: The tool lowers the barrier to entry for researchers new to computational glycobiology.
\end{enumerate}

%==============================================
% CONCLUSIONS
%==============================================
\section{Conclusions}

GlycoSMILES2BAP provides an automated solution to the stereochemistry preservation problem in AlphaFold 3 glycan modeling. By converting standard glycan not